\section{Conclusions and the boundary of the series}\label{sec:conclusions}
The author's paired upper/lower picture admits a precise conditional realization. With the median contacts fixed and normal crossing required, a unique circle in the selected median plane has a centre and radius determined by the patch geometry. Its contacts lie inside finite triangles, and the entire circle avoids the material panels. The six contacts cannot belong to one planar circle; three symmetry-related local circles are possible and mutually disjoint. The material object remains the inherited annulus. None of these statements fixes a physical loop count or tube thickness.

The complete quarter-turned observer placement has an explicit centre, pivot convention, frame, scale and translation. Full-panel clearance follows from enclosing balls. The H and R comparisons distinguish height from radius, retain 84 distinct sites among 96 labels, and expose the difference between mirrored centres and mirrored bodies. Their placement rings and the local circles remain separate geometries.

Under the first observation law, complete frames reconstruct a registered vector, but repetition adds no noiseless source dimension. For fixed known staged calibration the generic native information is $(I,\Ch)$, with real rank four and explicit invisible directions, singular strata and inversion limits. The cubic counterexample identifies data that those records omit. A conditional twelve-site spectral assignment does not eliminate that omission or the source contraction's tie preference.

Under the second law, declared delays recover the uniform baseline when they vanish. They supply exact timing and parity comparisons, a weighted transfer bound, and an aligned pure-delay component whose gain is independent of distance. Unequal phasor weights give an all-frequency inverse at aligned equatorial sites; equal axial weights permit complete blind frequencies elsewhere. The contact half-value and independent-channel storage calculation specify the model's limitations without invalidating its internal mathematics.

Subsequent papers address an explicitly driven shared carrier and exterior reception, followed by compact encodings, collisions and nonlinear saturation. Those models are distinct from the conditional geometric and observation maps established here. What remains open here is the adoption of a physical channel-to-space dictionary, a shared material connection, calibrated source work and reciprocal coupling. Those choices are not prerequisites for the geometric and observation theorems just proved.

\section*{Acknowledgements}
\addcontentsline{toc}{section}{Acknowledgements}
Hilmir Fr\'imann Halld\'orsson originated and directed the paired-gap, transverse-loop, twelve-module and multi-ring questions developed here. The work forms part of a sustained human--AI collaboration, including earlier exploration with GPT-4o, GPT-assisted formulation and internal mathematical review, and Codex-assisted source recovery, derivation, checks, figure preparation and typesetting. The internal reviews are not external peer review.

\appendix
\section{Exact coordinates and a reproducibility dictionary}\label{app:coords}
The panel maps \eqref{eq:panels}, local vertices \eqref{eq:octagon} and welded cycles in Section~\ref{sec:shell} define all coordinates without software. For direct reference, the eighteen global vertices are listed below. The notation $h_0=(\sqrt2-1)/2$ shortens the six seam endpoints. The ordering matches Paper C and the retained coordinate data.
\begin{center}\small
\begin{tabular}{@{}rccc@{}}\toprule
ID & $x$ & $y$ & $z$\\\midrule
0 & $0$ & $\sqrt3/2$ & $-h_0$\\
1 & $-1/2+\sqrt2/4$ & $\sqrt6/4$ & $-1/2$\\
2 & $-\sqrt2/4$ & $\sqrt3/2-\sqrt6/4$ & $-1/2$\\
3 & $-1/2$ & $0$ & $-h_0$\\
4 & $-1/2$ & $0$ & $h_0$\\
5 & $-\sqrt2/4$ & $\sqrt3/2-\sqrt6/4$ & $1/2$\\
6 & $-1/2+\sqrt2/4$ & $\sqrt6/4$ & $1/2$\\
7 & $0$ & $\sqrt3/2$ & $h_0$\\
8 & $-h_0$ & $0$ & $-1/2$\\
9 & $h_0$ & $0$ & $-1/2$\\
10 & $1/2$ & $0$ & $-h_0$\\
11 & $1/2$ & $0$ & $h_0$\\
12 & $h_0$ & $0$ & $1/2$\\
13 & $-h_0$ & $0$ & $1/2$\\
14 & $\sqrt2/4$ & $\sqrt3/2-\sqrt6/4$ & $-1/2$\\
15 & $1/2-\sqrt2/4$ & $\sqrt6/4$ & $-1/2$\\
16 & $1/2-\sqrt2/4$ & $\sqrt6/4$ & $1/2$\\
17 & $\sqrt2/4$ & $\sqrt3/2-\sqrt6/4$ & $1/2$\\\bottomrule
\end{tabular}
\end{center}
The three face cycles are $(0,1,2,3,4,5,6,7)$, $(3,8,9,10,11,12,13,4)$ and $(10,14,15,0,7,16,17,11)$. Each consecutive pair, including closure, is a material edge; the three shared pairs are counted only once. The measurement bases in Section~\ref{sec:shell} are not appended to these cycles.

The coordinate formulas above, together with the stated panel cycles, give a complete geometric specification. The accompanying source map identifies the retained coordinate data and the definitions used for each construction.

\section{Manuscript-local verification and figure construction}\label{app:checks}
The mathematical arguments in the main text establish the all-parameter statements. The manuscript checks verify displayed algebraic identities, rational inequality certificates, the fixed tables and the coordinates used in this edition. The static figure renderer uses retained coordinates, frames and selected state-contrast data. Circle and analytic phase samples are drawing primitives, not time steps or native trajectories.

The figure presentation separates solid material edges, dashed measurement closures, conditional local curves and reference axes. Figures~\ref{fig:geometry}--\ref{fig:delay} use the retained data, with presentation choices limited to camera, typography and panel layout. No new observer position, carrier or finite tube thickness is selected.

Finite checks support the transcription of the displayed formulas and data; the proofs establish their stated domains of validity. Counterexamples and exceptional cases delimit those domains.
